\section{Проектирование GPU-вычислительного модуля}

\subsection{Обоснование выбора Vulkan Compute}

Для реализации GPU-ускорения физического движка были рассмотрены четыре основных технологии: CUDA, OpenCL, Metal Compute и Vulkan Compute. Выбор технологии определяется требованиями к кроссплатформенности, интеграции с существующим рендерингом и возможности использования единого GPU-устройства для симуляции и визуализации.

CUDA обеспечивает наилучшую производительность и богатейшую экосистему инструментов (Thrust, cuBLAS, cuSPARSE), однако ограничена исключительно оборудованием NVIDIA, что нарушает требование кроссплатформенности~--- один из ключевых принципов архитектуры «Модейлера».

OpenCL является открытым стандартом с поддержкой широкого класса ускорителей, однако его развитие фактически заморожено: последняя значимая версия спецификации (OpenCL~3.0) вышла в 2020~г., и крупнейшие поставщики GPU de facto перенесли акцент на проприетарные или Vulkan-основанные решения~\cite{vulkan13spec}.

Metal~Compute доступен только на платформах Apple и поэтому не рассматривается.

\textbf{Vulkan Compute} выбран как основная технология GPU-ускорения. Его преимущества в контексте данной работы:
\begin{enumerate}
  \item \textit{Интеграция с рендерингом.} VulkanSceneGraph уже инициализирует Vulkan-устройство; вычислительные конвейеры разделяют с рендерингом логическое устройство (\texttt{vk::Device}), очереди и буферную память, исключая копирование через системную шину.
  \item \textit{Кроссплатформенность.} Vulkan поддерживается на Windows (NVIDIA, AMD, Intel), Linux и Android.
  \item \textit{Явное управление синхронизацией.} Примитивы Vulkan (fence, semaphore, pipeline barrier) позволяют точно задать порядок исполнения между вычислительными и графическими командами.
  \item \textit{Шейдеры в формате SPIR-V.} Единое промежуточное представление позволяет компилировать шейдеры из GLSL при сборке и распространять бинарные SPIR-V модули.
\end{enumerate}

\subsection{Организация вычислительного конвейера}

GPU-модуль состоит из трёх независимых вычислительных конвейеров (\textit{compute pipelines}), соответствующих трём стадиям симуляционного шага, выполняемым на GPU (рис.~\ref{fig:sim_pipeline}): широкой фазе, узкой фазе и решателю. Каждый конвейер включает один или несколько вычислительных шейдеров (GLSL compute shaders, скомпилированных в SPIR-V).

Рисунок~\ref{fig:gpu_pipeline} показывает граф вычислений для одного шага симуляции.

\begin{figure}[h]
\centering
\begin{tikzpicture}[
  font=\small,
  cpubox/.style={draw, rounded corners=3pt, fill=orange!15,
                 minimum width=5.5cm, minimum height=0.75cm, align=center},
  gpubox/.style={draw, rounded corners=3pt, fill=blue!15,
                 minimum width=5.5cm, minimum height=0.75cm, align=center},
  barrierbox/.style={draw, dashed, fill=gray!10,
                     minimum width=5.5cm, minimum height=0.5cm, align=center,
                     font=\footnotesize\itshape},
  arr/.style={-{Stealth[length=5pt]}, thick},
  label/.style={font=\footnotesize, text=gray},
]

% CPU side
\node[cpubox] (cpu_prep)   at (-3.2, 0)    {CPU: обновление SSBO (SoA данные)};
\node[cpubox] (cpu_sync)   at (-3.2,-6.6)  {CPU: чтение результатов SSBO};
\node[cpubox] (cpu_scene)  at (-3.2,-7.7)  {CPU: обновление трансформаций VSG};

% GPU side
\node[gpubox] (broad)      at (3.2,-1.1)   {Broadphase: AABB-пары};
\node[barrierbox] (bar1)   at (3.2,-2.0)   {pipeline barrier};
\node[gpubox] (narrow)     at (3.2,-3.1)   {Narrowphase: GJK/EPA};
\node[barrierbox] (bar2)   at (3.2,-4.0)   {pipeline barrier};
\node[gpubox] (solver)     at (3.2,-5.1)   {PGS-solver ($n$ итераций)};
\node[barrierbox] (bar3)   at (3.2,-6.0)   {pipeline barrier};

% Arrows CPU->GPU and back
\draw[arr] (cpu_prep.east)  -- ++(0.5,0) |- (broad.west)
    node[midway, above, label] {vkCmdCopyBuffer};
\draw[arr] (solver.south) -- ++(0,-0.2) -| (-3.2,-6.35)
    node[midway, below, label] {fence};
\draw[arr] (cpu_sync.south) -- (cpu_scene.north);

% Arrows within GPU
\draw[arr] (broad.south) -- (bar1.north);
\draw[arr] (bar1.south)  -- (narrow.north);
\draw[arr] (narrow.south)-- (bar2.north);
\draw[arr] (bar2.south)  -- (solver.north);
\draw[arr] (solver.south)-- (bar3.north);

% Labels
\node[label, left=0.3cm of cpu_prep]  {CPU};
\node[label, right=0.3cm of broad]    {GPU};

\end{tikzpicture}
\caption{Граф вычислений GPU-модуля физического движка на одном шаге симуляции}
\label{fig:gpu_pipeline}
\end{figure}

Все три GPU-стадии записываются в единый \texttt{VkCommandBuffer} перед отправкой. Барьеры памяти между ними задают тип доступа и стадию конвейера источника/получателя, чтобы гарантировать видимость результатов предыдущей стадии для следующей. Функция~\texttt{vkQueueSubmit} отправляет командный буфер в вычислительную очередь; CPU-поток ожидает завершения через~\texttt{VkFence}.

\subsection{Структуры буферов данных}
\label{sec:ssbo_layout}

Все данные физического движка, передаваемые в GPU, размещаются в \textit{Shader Storage Buffer Objects} (SSBO)~--- буферах, доступных для чтения и записи из вычислительных шейдеров. Использование нескольких специализированных SSBO вместо одного общего упрощает управление барьерами и сокращает объём передаваемых данных.

\textbf{BodiesSSBO} хранит состояние всех тел в SoA-формате (каждый атрибут --- отдельный массив). Размер буфера фиксируется при создании сцены на основе максимально ожидаемого числа тел $N_{\max}$. Структура буфера в GLSL:

\begin{lstlisting}[language=C, caption={Структура основного буфера тел (GLSL)}]
layout(std430, binding = 0) buffer BodiesSSBO {
    vec4  positions[N_MAX];      // xyz + padding
    vec4  orientations[N_MAX];   // кватернион xyzw
    vec4  linear_vels[N_MAX];    // xyz + inv_mass
    vec4  angular_vels[N_MAX];   // xyz + padding
    mat4  inv_inertia_world[N_MAX]; // 3x3 в mat4
    uint  flags[N_MAX];          // бит 0: static, бит 1: sleeping
};
\end{lstlisting}

Линейная скорость и обратная масса объединены в один \texttt{vec4} (паттерн \textit{struct packing}), сокращая число обращений к памяти на~25\%.

\textbf{BroadphaseSSBO} хранит оси AABB для всех тел и выходной список потенциально пересекающихся пар:

\begin{lstlisting}[language=C, caption={Буфер широкой фазы (GLSL)}]
layout(std430, binding = 1) buffer BroadphaseSSBO {
    vec4 aabb_min[N_MAX];        // минимальная точка AABB
    vec4 aabb_max[N_MAX];        // максимальная точка AABB
    uint pair_count;             // атомарный счётчик пар
    uint pairs[MAX_PAIRS * 2];   // чётные — bodyA, нечётные — bodyB
};
\end{lstlisting}

Атомарный счётчик~\texttt{pair\_count} обновляется через~\texttt{atomicAdd} в шейдере: каждый поток, обнаруживший перекрытие AABB, атомарно резервирует слот в массиве~\texttt{pairs}.

\textbf{ContactsSSBO} хранит результаты узкой фазы --- контактные точки:

\begin{lstlisting}[language=C, caption={Буфер контактных точек (GLSL)}]
struct Contact {
    vec4  normal_depth;   // нормаль xyz + глубина w
    vec4  point_a;        // точка контакта на теле A
    vec4  point_b;        // точка контакта на теле B
    uint  body_a;
    uint  body_b;
    uint  feature_id;     // для warm-starting
    float lambda_n;       // нормальный импульс (накопленный)
    vec2  lambda_t;       // тангенциальные импульсы
    uint  _pad[2];
};

layout(std430, binding = 2) buffer ContactsSSBO {
    uint    contact_count;
    Contact contacts[MAX_CONTACTS];
};
\end{lstlisting}

Поле \texttt{lambda\_n} и \texttt{lambda\_t} хранят накопленные импульсы из предыдущего шага и используются для warm-starting в начале работы решателя.

\subsection{Широкая фаза на GPU}

Вычислительный шейдер широкой фазы запускается с рабочей группой размером $64 \times 1 \times 1$; число групп $\lceil N^2 / 128 \rceil$ охватывает все пары тел. Шейдер выполняет наивную проверку AABB: каждый поток вычислительной группы обрабатывает одну пару $(i, j)$ с $i < j$:

\begin{lstlisting}[language=C, caption={Фрагмент шейдера широкой фазы}]
layout(local_size_x = 64) in;

void main() {
    uint pair_idx = gl_GlobalInvocationID.x;
    // Декодировать линейный индекс пары в (i, j)
    uint i = uint((-1.0 + sqrt(1.0 + 8.0*pair_idx)) / 2.0);
    uint j = pair_idx - i*(i+1)/2;
    if (i >= body_count || j >= body_count) return;

    // Пропустить статические пары
    if ((flags[i] & 1u) != 0 && (flags[j] & 1u) != 0) return;

    // Проверка перекрытия AABB
    bool overlap =
        all(lessThan(aabb_min[i].xyz, aabb_max[j].xyz)) &&
        all(lessThan(aabb_min[j].xyz, aabb_max[i].xyz));

    if (overlap) {
        uint slot = atomicAdd(pair_count, 1u);
        if (slot < MAX_PAIRS) {
            pairs[slot*2]   = i;
            pairs[slot*2+1] = j;
        }
    }
}
\end{lstlisting}

Наивный $O(N^2)$ алгоритм широкой фазы эффективен при числе тел $N \lesssim 1000$, поскольку GPU обрабатывает все пары параллельно за один запуск шейдера. При $N > 1000$ планируется переход к алгоритму, основанному на сортировке по одной из осей (\textit{Sort and Sweep}), реализованному через Vulkan-совместимый алгоритм параллельной сортировки Bitonic Sort.

\subsection{Узкая фаза на GPU}

Вычислительный шейдер узкой фазы обрабатывает каждую пару из~\texttt{BroadphaseSSBO} независимо: один поток --- одна пара. Реализован алгоритм GJK (описанный в разделе~2.2) в варианте, пригодном для исполнения на GPU: без рекурсии, с фиксированным числом итераций (не более 32), с выходом по достижению нулевого симплекса или предела итераций.

При обнаружении столкновения (нулевое минимальное расстояние) запускается алгоритм EPA для вычисления нормали и глубины проникновения. Результат записывается в~\texttt{ContactsSSBO} через тот же~\texttt{atomicAdd}-механизм.

Ключевое ограничение GPU-GJK: функция опоры для каждого примитива должна быть реализована без динамической диспетчеризации (виртуальных вызовов). Тип коллайдера закодирован полем~\texttt{uint collider\_type} в отдельном буфере, и шейдер использует~\texttt{switch}-оператор для выбора нужной функции опоры. Это типичная техника для GPU-кода, где виртуальные таблицы функций не поддерживаются.

\subsection{Решатель на GPU}

PGS-решатель выполняется итерационно: каждая итерация состоит из одного запуска вычислительного шейдера. Шейдер обрабатывает все контакты параллельно, применяя $\Delta\lambda$ к скоростям тел.

Основная проблема параллельного PGS --- \textit{гонки данных} при записи в скорости тел: несколько потоков могут одновременно читать и обновлять скорость одного тела. Используется следующая стратегия: тела разбиваются на независимые \textit{цветовые классы} (\textit{graph coloring}) --- такие, что в каждом классе нет двух тел, связанных одним контактом. Итерация решателя выполняется поочерёдно для каждого цветового класса; внутри одного класса все контакты обрабатываются параллельно без конфликтов.

Граф расцветки строится на CPU при обнаружении новых контактов и остаётся стабильным, пока топология контактов не изменяется. При типичных сценах число цветов не превышает $4$--$6$, что означает $4$--$6$ последовательных запусков шейдера на итерацию PGS~--- накладные расходы приемлемы при числе тел, достаточно большом для того, чтобы GPU-параллелизм давал выигрыш.

\subsection{Синхронизация CPU и GPU}

Схема синхронизации разработана с целью минимизации времени простоя как CPU, так и GPU. Используется двойная буферизация SSBO: пока GPU обрабатывает шаг $t$, CPU подготавливает данные для шага $t+1$ во втором наборе буферов. Переключение буферов происходит после завершения GPU-шага, подтверждённого~\texttt{VkFence}.

Временна́я диаграмма одного кадра симуляции показана на рисунке~\ref{fig:cpu_gpu_timeline}.

\begin{figure}[h]
\centering
\begin{tikzpicture}[font=\small, x=1.0cm, y=0.7cm]
  % CPU timeline
  \draw[thick] (0,2) -- (12,2);
  \node[left] at (0,2) {CPU};
  \fill[orange!50] (0,1.6) rectangle (2,2.4) node[midway, font=\footnotesize] {обн. SSBO};
  \fill[orange!50] (2.2,1.6) rectangle (3.5,2.4) node[midway, font=\footnotesize] {submit};
  \fill[gray!20]   (3.5,1.6) rectangle (7,2.4) node[midway, font=\footnotesize] {ожидание fence};
  \fill[green!40]  (7,1.6) rectangle (9.5,2.4) node[midway, font=\footnotesize] {чтение SSBO};
  \fill[green!40]  (9.5,1.6) rectangle (11.5,2.4) node[midway, font=\footnotesize] {обн. VSG};

  % GPU timeline
  \draw[thick] (0,0) -- (12,0);
  \node[left] at (0,0) {GPU};
  \fill[gray!20]  (0,-0.4) rectangle (3.5,0.4) node[midway, font=\footnotesize] {простой};
  \fill[blue!40]  (3.5,-0.4) rectangle (5,0.4) node[midway, font=\footnotesize] {broadphase};
  \fill[blue!60]  (5,-0.4) rectangle (7,0.4) node[midway, font=\footnotesize] {narrowphase};
  \fill[blue!80]  (7,-0.4) rectangle (9.2,0.4) node[midway, font=\footnotesize, text=white] {PGS solver};
  \fill[gray!20]  (9.2,-0.4) rectangle (12,0.4) node[midway, font=\footnotesize] {простой};

  % Fence arrow
  \draw[dashed, -{Stealth[length=5pt]}] (9.2,0.4) -- (7,1.6)
    node[midway, above, sloped, font=\footnotesize] {fence signal};

  % Time axis
  \draw[->] (0,-1) -- (12,-1) node[right] {$t$};
  \node[below] at (0,-1) {$t_0$};
  \node[below] at (12,-1) {$t_0 + \Delta t$};
\end{tikzpicture}
\caption{Временна́я диаграмма взаимодействия CPU и GPU на одном шаге симуляции}
\label{fig:cpu_gpu_timeline}
\end{figure}

Основной источник задержки --- время ожидания CPU завершения GPU-шага. При достаточно большом числе тел время GPU-вычислений превышает время CPU-подготовки, и ожидание становится пренебрежимо малым относительно общего шага. При малом числе тел ($N \lesssim 100$) накладные расходы на запуск шейдеров и барьеры превышают выигрыш от параллелизма; в этом случае движок автоматически переключается на CPU-реализацию.

\subsection{Интеграция с Vulkan-устройством «Модейлера»}

Физический движок использует то же~\texttt{vk::Device}, что и слой рендеринга VSG, совместно с ним управляя памятью через~\texttt{vmaCreateBuffer} (Vulkan Memory Allocator). Вычислительные конвейеры создаются отдельно от графических, используя тип~\texttt{VK\_STRUCTURE\_TYPE\_COMPUTE\_PIPELINE\_CREATE\_INFO}.

Для вычислительных команд используется отдельная вычислительная очередь (compute queue family), что позволяет GPU исполнять физику и рендеринг асинхронно при наличии аппаратной поддержки нескольких очередей. Синхронизация между рендерингом и физикой обеспечивается~\texttt{VkSemaphore}: рендеринг нового кадра начинается только после того, как физический шаг записал обновлённые трансформации в буфер, доступный вершинному шейдеру.
